\section{预期贡献总结}

本文的预期贡献涵盖规程设计、系统设计、方法学与实证发现四个层面。
贡献 0：规程设计。提出 Pilot→PreScreen→Main 的三段式可审计规程，在进入主实验前预注册并锁定阈值、超参数与权重映射。流程允许在软件未完全自动化的情况下通过任务包分发执行，但必须保留可审计的样本隔离与分发清单。
贡献 1：系统设计。提出可审计的半自动标注体系。三级报告口径 T/I/M 避免选择偏差；所有门控失败与 OOS 可追溯；Meta-label 作为决策副产品提供分布监控信号。
贡献 2：方法学。预筛选测试样本、场景特异可靠度、OOD 触发路由的组合。$r_{u}^{(s)}$ 的估计协议与 LOO 原则；Per-tag 共识与分歧度计算；OOD 风险分数 $d_{t}$ 的预注册主实现与敏感性分析框架。
贡献 3：实证发现。在全景布局标注任务上验证，以最终实证结果为准。半自动初始化在相同图像对照下可降低 \texttt{active\_time}；场景与分布感知路由在相同预算下可提升一致性代理并改善预算效率；门控失败与字段缺失等作为安全性/数据可用性审计事件被完整追踪与披露（不作为主收益主张来源）；反例案例库揭示三类主要失效模式，并形成过程性证据链。
对标文献：Kara et al.（JAIR 2018）聚焦主动学习与共识估计，本文在此基础上引入场景分层；Liu et al.（InfoVis 2018）聚焦 worker 分型与交互验证，本文采用离线路由模拟加以拓展。创新点在于将可审计性、场景感知与元标签共识进行系统性整合。


